\section{Introdução e Problema da Inferência}

\begin{frame}{Como Conhecer a População?}
    \begin{block}{O Problema da Inferência}
        \begin{itemize}
            \item Temos acesso apenas a uma \textbf{amostra} de tamanho $N$.
            \item Nosso objetivo é descrever ou estimar propriedades da \textbf{população} (como a média $\mu$).
            \item Não temos orçamento ou meios para auditar a população completa.
        \end{itemize}
    \end{block}
    
    \begin{alertblock}{As Duas Abordagens Principais}
        \begin{enumerate}
            \item \textbf{Abordagem Clássica:} Utiliza fórmulas matemáticas baseadas em suposições rígidas de distribuição (ex: Teorema Central do Limite).
            \item \textbf{Abordagem Computacional:} Utiliza simulação e reamostragem a partir dos próprios dados observados (Bootstrap).
        \end{enumerate}
    \end{alertblock}
\end{frame}

\begin{frame}{O que é um Teste A/B?}
    \begin{block}{Definição Básica}
        Um experimento controlado para comparar duas versões de uma variável (A e B) com o objetivo de determinar qual delas é mais eficaz.
    \end{block}
    
    \begin{columns}
        \begin{column}{0.5\textwidth}
            \textbf{Grupo A (Controle)}
            \begin{itemize}
                \item Versão clássica do sistema.
                \item Experiência do usuário padrão.
                \item Medicamento placebo (em saúde).
            \end{itemize}
        \end{column}
        \begin{column}{0.5\textwidth}
            \textbf{Grupo B (Tratamento)}
            \begin{itemize}
                \item Nova funcionalidade ou design.
                \item Novo layout de checkout rápido.
                \item Novo princípio ativo.
            \end{itemize}
        \end{column}
    \end{columns}
    
    \vspace{0.4cm}
    \centering
    \textit{Queremos provar que a diferença observada entre os grupos é real, e não mera flutuação aleatória!}
\end{frame}

\begin{frame}{Exemplo Matemático: Estimativa de Parâmetros}
    \begin{block}{Problema de Negócio}
        Imagine que medimos o tempo de carregamento de uma página em 5 visitas (em segundos):
        $$X = \{1.5, 2.3, 1.8, 3.0, 2.4\}$$
        A estimativa pontual da média populacional $\mu$ é a média amostral $\bar{X}$:
        $$\bar{X} = \frac{1.5 + 2.3 + 1.8 + 3.0 + 2.4}{5} = \frac{11.0}{5} = 2.2\text{ s}$$
    \end{block}
    \begin{block}{Questão Fundamental}
        Essa estimativa pontual ($2.2\text{ s}$) é representativa? Se coletarmos outra amostra de 5 visitas, qual será a variação da média?
    \end{block}
\end{frame}

\begin{frame}[fragile]{Demonstração em Python: Média Amostral e Variabilidade}
    \begin{lstlisting}[language=Python]
import numpy as np

# Amostra de tempos de carregamento (em segundos)
X = np.array([1.5, 2.3, 1.8, 3.0, 2.4])

# Estimativa pontual da media
media_amostral = np.mean(X)
print(f"Media amostral: {media_amostral:.2f} s")
# Saida: Media amostral: 2.20 s

# Simulando a variabilidade amostral (se a populacao fosse conhecida)
np.random.seed(42)
populacao = np.random.normal(loc=2.2, scale=0.5, size=10000)
novas_medias = [np.random.choice(populacao, size=5).mean() for _ in range(1000)]
print(f"Desvio padrao das medias (erro padrao): {np.std(novas_medias):.3f}")
    \end{lstlisting}
\end{frame}

